\section{Related Work}
\label{appendix_sec:realted_work}


% \paragraph{Fast weight programming and test-time-training.} Test-time-training can be viewed as a special case of fast weight programming where the update rule is derived from the gradient of an online learning objective.  In the context of sequence modeling, TTT is essentially a \emph{neural memory} trying to encode KV cache into the neural memory by optimizing the regression loss of reconstructing values based on keys with a neural encoder whose weight is the RNN's hidden state.

% and in the context of neural sequence modeling, it is usually a l2-based regression loss from the keys and values with proper regularization such as l2norm-based (recurrent) weight matrix's norm (cf. Test-Time Regression \cite{wang2025testtimeregressionunifyingframework} work for a more holistic summarization).  .but could also be extended to other objective such as with $l_p$ norm \cite{behrouz2025itsconnectedjourneytesttime}.

\paragraph{Test-time training.} Test-Time Training (TTT)~\cite{sun2024learning}  is an emerging concept in sequence modeling that extends the concept of recurrent states in RNNs to online-adapted neural network components. In TTT models, a subset of weights, termed "fast weights," are updated to learn in-context. Existing methods typically employ a self-supervised loss that encourages these fast weights to memorize key-value associations from in-context tokens, using variants of gradient descent for online adaptation. TTT~\cite{sun2024learning,wang2025testtimeregressionunifyingframework} has opened a vast design space for new recurrent model architectures. For instance, many recent works have developed novel test-time optimizers~\cite{behrouz2024titans, karami2025lattice} and online training objectives~\cite{behrouz2025s}. However, current TTT approaches often suffer from low hardware utilization and limited state sizes, and consequently have not yet demonstrated their full potential.  Our work primarily addresses these challenges by advocating for a new paradigm of using extremely large online minibatch (chunk) sizes for updating the fast weights. This paradigm can achieve orders-of-magnitude higher hardware utilization without relying on error-prone custom kernel implementations. Furthermore, it enables efficient scaling of nonlinear state sizes and offers the flexibility to use diverse fast weight neural networks and optimizers, thereby accelerating research progress in this area.


\paragraph{Combining chunk attention with recurrence.}
Several recent models combine local chunk attention with linear recurrence, such as Gated Attention Unit (GAU) \cite{hua2022transformerqualitylineartime}, MEGA \cite{ma2023mega}, MEGALODON \cite{ma2024megalodon}, and InfiniAttention \cite{munkhdalai2024leavecontextbehindefficient}. Among these, InfiniAttention is conceptually closest to our work, as it incorporates recurrence at the chunk level using the delta rule—interpreted as an online linear regression objective from the perspective of Test-Time Training (TTT). However, this update rule is limited in expressivity. In contrast, we employ a significantly more expressive update mechanism derived from a more general TTT framework, and demonstrate the substantial gains this brings.

Block-Recurrent Transformer \cite{hutchins2022blockrecurrent} also explores large chunk memory updates, where memory tokens act as recurrent states that can self-attend and cross-attend with input tokens during each chunk update via attention mechanisms.  The Perceiver-style register-token attention baseline used in our novel view synthesis experiments (Sec.~\ref{sec:exp_nvs}, Table~\ref{tab:attention_comparison}) is conceptually similar to the Block-Recurrent Transformer in its use of register tokens for context compression. As shown in Figure~\ref{fig:view_syn_results}, our method significantly outperforms this approach in both speed and quality, with a comparable state size.

%However, its update mechanism is largely heuristic and lacks a principled foundation. Our approach, instead, adopts a well-grounded update rule rooted in TTT formalism.


\paragraph{Novel view synthesis.} Novel view synthesis (NVS) is a long-standing task at the intersection of computer vision, graphics, and computational photography, requiring algorithms to render images of a static scene from previously unobserved viewpoints. Optimization-based approaches, such as NeRF~\cite{mildenhall2021nerf} and 3D Gaussian Splatting~\cite{kerbl20233d}, have achieved significant breakthroughs.
These methods optimize a set of parameterized graphics primitives (i.e., explicit or implicit representations of radiance fields) through differentiable volumetric rendering to minimize reconstruction loss on input images. After an optimization process typically lasting tens of minutes, these approaches can render novel views photorealistically, and the optimized parameters form a 3D representation of the input scene.

Recently, data-driven approaches~\cite{zhang2024gs, jin2024lvsm, ziwen2024long, liu2023zero} have also shown promising results. These methods can either directly render novel views or predict 3D representations given input images. Although successful on simpler object datasets, these methods often struggle with densely sampled scenes (e.g., scenes with over 100 input images).  Our experiments demonstrate that our large-chunk test-time training approach outperforms or achieves comparable performance to 3D Gaussian Splatting on challenging scene datasets with up to 128 input images with $960\times536$ resolution at challenging scene datasets.We hope our method will inspire further research into effectively scaling data-driven NVS methods to longer and more complex input sequences.



\paragraph{Autoregressive video diffusion.}
Current state-of-the-art video generation is dominated by bidirectional diffusion transformers operating in latent space~\cite{videoworldsimulators2024, yang2024cogvideox, polyak2410movie, wang2025wan}. These methods factorize the video distribution into a sequence of conditional distributions based on noise levels, following diffusion processes~\cite{sohl2015deep, song2020score} or flow matching~\cite{lipman2022flow}, then use a diffusion transformer to jointly learn all the conditional distribution. Autoregressive video diffusion~\cite{alonso2024diffusion, jin2024pyramidal, valevski2024diffusion, ruhe2024rolling, yin2024slow, song2025history} introduces an additional temporal dimension to this factorization, where the neural networks learns to model the conditional probability of the next chunks of videos at different noise levels, conditional on previous videos and noisier  version of current video frames.  

During training, some autoregressive methods employ teacher forcing, supervising the model on noisy video frames given previous  clean context frames as condition~\cite{alonso2024diffusion,jin2024pyramidal,valevski2024diffusion}, though this can lead to low token utilization, i.e. only a small portion of tokens get supervision. To improve token efficiency, other techniques such as progressive noise injection~\cite{ruhe2024rolling} or the use of frame-independent noises (sometimes in a diffusion-forcing style)~\cite{yin2024slow, chen2024diffusion, magi1} have been proposed. When applying our large-chunk design to autoregressive video generation, we format the input sequence with interleaved clean and noisy chunks (see Equation~\ref{eq:video_sequence_basic}). This strategy achieves over 50\% token utilization and integrates effectively with our large-chunk TTT implementation, by only changing a few lines to constrain fast-weights are only updated on clean frame chunks.